\documentclass[11pt]{article}
\usepackage[margin=1in]{geometry}
\usepackage{lmodern}
\usepackage[T1]{fontenc}
\usepackage{microtype}
\usepackage{amsmath,amssymb,mathtools}
\usepackage{booktabs}
\usepackage{enumitem}
\usepackage{xcolor}
\usepackage[most]{tcolorbox}
\usepackage[colorlinks=true,linkcolor=blue!50!black,urlcolor=blue!50!black]{hyperref}

\definecolor{keyblue}{RGB}{30,70,140}
\newtcolorbox{keyidea}[1][]{colback=keyblue!5,colframe=keyblue!70,
  fonttitle=\bfseries,title=#1,boxrule=0.6pt,arc=2pt,left=6pt,right=6pt}
\newtcolorbox{caution}[1][]{colback=orange!6,colframe=orange!70!black,
  fonttitle=\bfseries,title=#1,boxrule=0.6pt,arc=2pt,left=6pt,right=6pt}

\newcommand{\E}{\mathbb{E}}
\newcommand{\mb}{\bar m}
\setlist{itemsep=2pt,topsep=4pt}

\title{\textbf{The HJB and Fokker--Planck Equations of Problem 21}\\[4pt]
\large AMS 516 Project --- step-by-step derivations and what they mean for the solvers}
\author{Deepak Singh}
\date{October 2026}

\begin{document}
\maketitle

\begin{abstract}
These notes derive the two equations of the mean-field portfolio game with a shortfall
penalty: the Hamilton--Jacobi--Bellman (HJB) equation for one investor's value function,
solved backward in time, and the Fokker--Planck (FP) equation for the population's wealth
distribution, solved forward in time. Each derivation goes one step at a time, is checked
against the $\lambda=0$ closed forms, and ends with what it implies for the finite-difference,
Monte Carlo and PINN solvers. Appendix~\ref{app:fd-hjb} turns the HJB into the
finite-difference scheme used in the code. Numbers use \texttt{configs/base.yaml}:
$r=0.02$, $\mu=0.08$, $\sigma=0.20$, $\gamma=1$, $\lambda=2$, $T=1$,
$m_0=\mathcal N(1,0.2^2)$, domain $x\in[-1,3]$.
\end{abstract}

\setcounter{tocdepth}{2}
\tableofcontents

%======================================================================
\section{The model in one place}

Each investor holds wealth $X_t$ and puts a dollar amount $\pi_t$ in the stock:
\[
  dX_t = \bigl[rX_t + \pi_t(\mu-r)\bigr]\,dt + \pi_t\sigma\,dW_t ,
\]
with their own Brownian motion $W_t$ (independent noise). They maximise
$\E[U(X_T,\mb_T)]$ with
\[
  U(x,\mb) = -e^{-\gamma x} - \tfrac{\lambda}{2}\bigl((\mb-x)^+\bigr)^2 ,
\]
where $\mb_T$ is the population's average terminal wealth. Throughout,
$\kappa=(\mu-r)/\sigma=0.3$ is the Sharpe ratio (\texttt{Params.kappa}) and the Merton
position is $(\mu-r)/(\gamma\sigma^2)=1.5$.

%======================================================================
\section{The HJB equation}

\subsection*{Step 1: The value function}
\addcontentsline{toc}{subsection}{Step 1: The value function}
\[
  V(t,x) = \sup_{\pi}\ \E\bigl[\,U(X_T,\mb_T)\ \big|\ X_t = x\,\bigr] .
\]
$V(t,x)$ is the best expected utility still available to an investor who has wealth $x$ at
time $t$ and invests optimally until $T$.
\begin{itemize}
  \item It depends only on the current time and wealth.
  \item $\mb_T$ is a \textbf{fixed number} here: one investor can't move the average.
  \item At $T$ there is nothing left to choose: $V(T,x)=U(x,\mb_T)$. This is the only place
        the utility enters.
\end{itemize}

\subsection*{Step 2: Dynamic programming principle}
\addcontentsline{toc}{subsection}{Step 2: Dynamic programming principle}
Split the remaining time into the next instant $dt$ and everything after:
\[
  V(t,x) = \sup_{\pi}\ \E\bigl[\,V(t+dt,\,X_{t+dt})\ \big|\ X_t = x\,\bigr] .
\]
Choose $\pi$ for the next instant; after that, act optimally, which is worth
$V(t+dt,X_{t+dt})$ by definition. There is no consumption, so nothing is collected along the
way.

\subsection*{Step 3: It\^o's formula}
\addcontentsline{toc}{subsection}{Step 3: It\^o's formula}
With $dX = [rx+\pi(\mu-r)]\,dt + \pi\sigma\,dW$ and $(dX)^2=\pi^2\sigma^2\,dt$,
\[
  dV = \Bigl(V_t + \bigl[rx+\pi(\mu-r)\bigr]V_x + \tfrac12\pi^2\sigma^2V_{xx}\Bigr)dt
       + \pi\sigma V_x\,dW .
\]
The $dW$ term has mean zero, so $\E[dV]$ is the bracket times $dt$.

\subsection*{Step 4: The HJB in sup form}
\addcontentsline{toc}{subsection}{Step 4: The HJB in sup form}
Substituting $V(t+dt,X_{t+dt})=V(t,x)+dV$ into the DPP, $V(t,x)$ cancels, leaving
$0=\sup_\pi\E[dV]$. Dividing by $dt$:
\[
  \boxed{\;V_t + \sup_{\pi}\Bigl\{\bigl[rx+\pi(\mu-r)\bigr]V_x + \tfrac12\pi^2\sigma^2V_{xx}\Bigr\} = 0,
  \qquad V(T,x)=U(x,\mb_T).\;}
\]
\begin{center}
\begin{tabular}{@{}ll@{}}
\toprule
Term & Meaning \\
\midrule
$V_t$ & how prospects change as time passes \\
$rx\,V_x$ & interest on wealth, valued at the marginal value of a dollar \\
$\pi(\mu-r)\,V_x$ & expected excess return from the stock \\
$\tfrac12\pi^2\sigma^2\,V_{xx}$ & cost of risk: variance times curvature (negative) \\
\bottomrule
\end{tabular}
\end{center}
The braces hold the same ``gain minus risk cost'' trade-off as the Merton problem,
with $V_x$ and $V_{xx}$ in place of the utility's slope and curvature.

The equation contains no utility: it depends only on the wealth dynamics. Plain
CARA, the original utility and the shortfall utility share it, and differ only in the
terminal condition.

\subsection*{Step 5: Choosing the best \texorpdfstring{$\pi$}{pi}}
\addcontentsline{toc}{subsection}{Step 5: Choosing the best position}
At fixed $(t,x)$, only the part depending on $\pi$ matters:
\[
  f(\pi) = \pi(\mu-r)V_x + \tfrac12\pi^2\sigma^2V_{xx},
\]
a parabola with $f''(\pi)=\sigma^2V_{xx}$.

\textbf{Second-order condition.} If $V_{xx}<0$ the parabola opens downward and has a single
maximum. If $V_{xx}\ge0$ the supremum is infinite (unlimited risk). Concavity of $V$ is
therefore essential (Step 9d shows it holds).

\textbf{First-order condition.} $f'(\pi)=0$ gives
\[
  \boxed{\;\pi^*(t,x) = -\frac{(\mu-r)\,V_x}{\sigma^2\,V_{xx}}
          = \frac{\mu-r}{\sigma^2}\cdot\frac{1}{A_V(t,x)},
  \qquad A_V = -\frac{V_{xx}}{V_x}.\;}
\]
$A_V$ is the \emph{indirect risk aversion}: how risk-averse the investor effectively is at
$(t,x)$, given optimal investing afterwards. At $t=T$, $V=U$ and $A_V=A(x)=-U''(x)/U'(x)$, the
risk aversion of the utility itself, so the terminal position is
$\pi^*(T,x)=1.5\cdot\gamma/A(x)$.

\textbf{Worked example at $t=T$, $\mb=1.1$.} Here
$f(\pi)=0.06\,U'\pi-\tfrac12\,0.04\,(-U'')\pi^2$.
\begin{itemize}
  \item Just above $\mb$: $U'\approx0.333$, $-U''\approx0.333$, so
        $f=0.0200\pi-0.00667\pi^2$, maximised at $\pi=1.5$.
  \item Just below $\mb$: $U'\approx0.333$ (continuous), $-U''\approx2.333$, so
        $f=0.0200\pi-0.0467\pi^2$, maximised at $\pi\approx0.214$.
\end{itemize}
The gain slope is the same on both sides; only the risk cost changes. The parabola below
$\mb$ is 7 times steeper, so its peak is 7 times closer to zero.

\textbf{Feedback rule.} $\pi^*$ depends on the state $(t,x)$ through $V_x$ and $V_{xx}$.
A solver learns $V$, and $\pi^*$ is read off its derivatives, so the derivatives must be
accurate, not just $V$.

\subsection*{Step 6: The reduced PDE}
\addcontentsline{toc}{subsection}{Step 6: The reduced PDE}
Substituting $\pi^*$,
\[
  \text{gain} = -\frac{(\mu-r)^2V_x^2}{\sigma^2V_{xx}} > 0,\qquad
  \text{risk cost} = \frac{(\mu-r)^2V_x^2}{2\sigma^2V_{xx}} = -\tfrac12\,\text{gain},
\]
so at the optimum the risk cost eats exactly half the expected gain, and
\[
  \boxed{\;V_t + rx\,V_x - \tfrac12\kappa^2\,\frac{V_x^2}{V_{xx}} = 0,
  \qquad V(T,x)=U(x,\mb_T).\;}
\]
With our numbers $\tfrac12\kappa^2=0.045$.

\textbf{Nonlinear.} The ratio $V_x^2/V_{xx}$ makes the equation nonlinear. For a fixed
strategy $\pi(t,x)$ it would be linear:
$V_t + [rx+\pi(\mu-r)]V_x + \tfrac12\pi^2\sigma^2V_{xx}=0$. The nonlinearity comes entirely
from the investor choosing the risk.

\textbf{Two invariances.} If $V$ solves the PDE, so do $V+c$ and $cV$ for $c>0$. The second
is why the original utility $-e^{-\gamma(x-\theta\mb)}$ changes nothing: its terminal
condition is a positive multiple of CARA's, so its solution is too, and the constant cancels
in $\pi^*$. The shortfall penalty is neither a shift nor a multiple.

\begin{caution}[Division by the second derivative]
If $V_{xx}$ is close to zero or positive anywhere, the residual blows up. Ways around it:
\begin{center}
\small
\begin{tabular}{@{}lll@{}}
\toprule
Approach & Idea & Watch out for \\
\midrule
Reduced PDE & train or solve on the residual above & keep $V$ concave \\
Multiplied through & $V_{xx}(V_t+rxV_x)-\tfrac12\kappa^2V_x^2=0$ & admits wrong solutions (constants) \\
Sup form + policy & learn $V$ and $\pi$ together & two unknowns, harder training \\
Policy iteration & fix $\pi$, solve the linear PDE, update $\pi$ & a few outer iterations \\
\bottomrule
\end{tabular}
\end{center}
\end{caution}

\subsection*{Step 7: Terminal and boundary conditions}
\addcontentsline{toc}{subsection}{Step 7: Terminal and boundary conditions}

\textbf{Terminal condition.} $V(T,x)=U(x,\mb_T)$, solved backward from $T$. Two features:
\begin{enumerate}
  \item It has a kink at $\mb_T$: $V$ and $V_x$ are continuous, $V_{xx}$ jumps by $\lambda$.
        Close to $T$, $V_{xx}$ and therefore $\pi^*$ change sharply near $\mb_T$.
  \item It contains $\mb_T$, which depends on everyone's $\pi^*$. This is the coupling to the
        Fokker--Planck equation (Section~\ref{sec:system}).
\end{enumerate}

\textbf{Why wealth boundaries.} The true problem lives on all of $\mathbb R$; the solver uses
$[-1,3]$. The cut is artificial, and what $V$ does at the edges is a modeling choice. A
condition is needed at both ends, because the diffusion $\pi^{*2}\sigma^2>0$ carries
information from both edges inward.

\textbf{Upper boundary, $x=3$.} The penalty can only matter if wealth falls from 3 below 1.1
by $T$, a drop of about 6 standard deviations ($\pi\sigma\approx0.3$ per $\sqrt{\text{year}}$).
So the investor is effectively CARA, and the closed form serves as a Dirichlet condition:
\[
  V(t,3) = -\exp\!\bigl(-3\gamma e^{r(T-t)} - \tfrac12\kappa^2(T-t)\bigr).
\]

\textbf{Lower boundary, $x=-1$.} Here both the exponential and the penalty matter and there is
no closed form. The options:
\begin{itemize}
  \item \textbf{Value of a fixed strategy (recommended).} If the investor holds the Merton
        position from $t$ on, $X_T$ is Gaussian with
        \[
          M = x\,e^{r(T-t)} + \frac{\kappa^2}{\gamma}(T-t),\qquad
          S^2 = \frac{\kappa^2}{\gamma^2}(T-t),\qquad d=\mb_T-M,
        \]
        and both parts of the expected utility have closed forms:
        \[
          V_{\text{fix}}(t,x) = -e^{-\gamma M+\gamma^2S^2/2}
          - \tfrac{\lambda}{2}\Bigl[(d^2+S^2)\,\Phi\!\bigl(\tfrac{d}{S}\bigr)
                                   + dS\,\varphi\!\bigl(\tfrac{d}{S}\bigr)\Bigr],
        \]
        with $\Phi,\varphi$ the standard normal CDF and density. Its CARA part reproduces the
        Merton value function exactly. Because the strategy isn't optimal near $x=-1$
        (optimal is about 2.2, Merton 1.5), this is a \emph{lower bound} on $V$. (Checked
        against Monte Carlo with $10^6$ paths: $-6.863$ vs $-6.861$ at $(t,x)=(0,-1)$.)
  \item \textbf{Push $x_{\min}$ far down and use CARA.} As $x\to-\infty$, $e^{-\gamma x}$
        outgrows the quadratic penalty and $A(x)\to\gamma$, but only for $x\lesssim-8$, where
        $V\sim e^{8}\approx3000$: a poor scale for a network.
  \item \textbf{Not} the option-pricing ``linearity'' condition $V_{xx}=0$: it makes
        $\pi^*$ infinite.
\end{itemize}

\textbf{The check that matters.} The population lives roughly in $[0,2.2]$ (at $T$, mean
$\approx1.11$, standard deviation $\approx0.36$). Move the boundary ($x_{\min}=-1\to-2$) and
confirm $V$ and $\pi^*$ don't change there.

\subsection*{Step 8: The CARA case solved exactly}
\addcontentsline{toc}{subsection}{Step 8: The CARA case solved exactly}
For $\lambda=0$, guess the shape of the terminal condition:
\[
  V(t,x) = -\exp\bigl(-\gamma a(t)x - b(t)\bigr),\qquad a(T)=1,\ b(T)=0 .
\]
Every derivative is $V$ times something simple:
\[
  V_t = (-\gamma a'x - b')V,\qquad V_x = -\gamma aV,\qquad V_{xx}=\gamma^2a^2V,
  \qquad \frac{V_x^2}{V_{xx}} = V .
\]
Substituting into the PDE and dividing by $V$:
\[
  -\gamma(a'+ra)\,x - \bigl(b'+\tfrac12\kappa^2\bigr) = 0\qquad\text{for all }x .
\]
Both the coefficient of $x$ and the constant must vanish, which gives two ODEs:
\[
  a' = -ra \;\Rightarrow\; a(t)=e^{r(T-t)},\qquad
  b' = -\tfrac12\kappa^2 \;\Rightarrow\; b(t)=\tfrac12\kappa^2(T-t),
\]
\[
  V(t,x) = -\exp\!\bigl(-\gamma x\,e^{r(T-t)} - \tfrac12\kappa^2(T-t)\bigr)
  \qquad(\texttt{merton\_value}).
\]
Then $A_V=\gamma e^{r(T-t)}$, independent of $x$, and $\pi^*=1.5\,e^{-r(T-t)}$: the Merton
position. The \emph{certainty equivalent}
$\mathrm{CE}=-\tfrac1\gamma\ln(-V)=x\,e^{r(T-t)}+\kappa^2(T-t)/(2\gamma)$ shows that access
to the stock is worth $\kappa^2/(2\gamma)=0.045$ of terminal wealth over one year. For the
original utility, $b(T)=-\gamma\theta\mb$, so $V_{\text{orig}}=e^{\gamma\theta\mb}V_{\text{CARA}}$.

The guess works because the PDE maps ``exponential of a linear function of $x$'' to itself,
collapsing into two ODEs.

\subsection*{Step 9: What the shortfall penalty changes}
\addcontentsline{toc}{subsection}{Step 9: What the shortfall penalty changes}
\begin{enumerate}[label=(\alph*)]
  \item \textbf{No closed form.} The terminal condition is not an exponential of a linear
        function, so the guess of Step 8 fails. $V$ must be computed numerically.
  \item \textbf{The position depends on wealth and time.} $A_V(t,x)$ now varies with $x$.
        At $T$ it is $A(x)$: 7.0 just below $\mb$, back to $\gamma$ at a gap of
        $1/\gamma$, and below $\gamma$ further down; for $t<T$ diffusion smooths the spike
        into a bump of width roughly $\pi\sigma\sqrt{T-t}\approx0.3\sqrt{T-t}$. Close to $T$
        the reaction near $\mb$ is sharp; early on it is milder and spread out.
  \item \textbf{The ``fear zone'' extends above $\mb$ before $T$.} An investor slightly above
        $\mb$ at $t<T$ may still finish below it, so their $V$ carries some expected penalty,
        $A_V>\gamma$, and they hold a little less than Merton.
  \item \textbf{Concavity is preserved.} Wealth is linear in $(x,\pi)$: averaging the
        strategies from two starting wealths gives the average wealth path. With a concave
        utility this makes $V(t,\cdot)$ concave for every $t$, so $V_{xx}<0$ and the reduced
        PDE of Step 6 is legitimate.
  \item \textbf{$V$ depends on $\mb_T$.} Strictly $V=V(t,x;\mb_T)$: moving $\mb_T$ moves the
        kink and the fear zone. Finding $\mb_T$ requires the population's wealth distribution,
        which is the Fokker--Planck equation.
\end{enumerate}

\begin{keyidea}[Where the HJB is hard]
Close to $T$ and near $x=\mb_T$: $V_{xx}$ changes fastest there and $\pi^*$ drops from 1.5 to
0.214. Finite-difference grids should be finest there; PINNs should sample extra collocation
points there.
\end{keyidea}

%======================================================================
\section{The Fokker--Planck equation}

\subsection*{Step 1: The wealth density}
\addcontentsline{toc}{subsection}{Step 1: The wealth density}
$m(t,x)$ is the density of investors at wealth $x$ at time $t$:
\[
  \text{fraction in }[a,b] = \int_a^b m(t,x)\,dx,\qquad \int m\,dx = 1,\qquad
  \mb(t) = \int x\,m(t,x)\,dx .
\]
It is both the smoothed histogram of the population and the probability density of one
typical investor's $X_t$; with independent noise and a continuum of investors these coincide
by the law of large numbers. Initially $m(0,\cdot)=\mathcal N(1,0.2^2)$.

\subsection*{Step 2: Derivation}
\addcontentsline{toc}{subsection}{Step 2: Derivation}
Everyone follows $\pi^*(t,x)$, which is now a \emph{known} function, not a choice. Name the
pieces of the dynamics:
\[
  b(t,x) = rx + \pi^*(t,x)(\mu-r),\qquad s^2(t,x) = \pi^*(t,x)^2\sigma^2,\qquad
  dX = b\,dt + s\,dW .
\]
\textbf{2a. It\^o on a test function.} For any smooth $\varphi(x)$ that vanishes, with its
derivative, at the edges:
$d\varphi(X) = [b\varphi' + \tfrac12s^2\varphi'']\,dt + s\varphi'\,dW$, so
\[
  \frac{d}{dt}\E[\varphi(X_t)] = \E\bigl[b\varphi' + \tfrac12s^2\varphi''\bigr].
\]
\textbf{2b. Write expectations as integrals against $m$:}
\[
  \frac{d}{dt}\int\varphi\,m\,dx = \int\bigl(b\varphi' + \tfrac12s^2\varphi''\bigr)m\,dx .
\]
\textbf{2c. Integrate by parts} (once on the first term, twice on the second; boundary
terms vanish):
\[
  \int\varphi\,\Bigl[m_t + \partial_x(bm) - \tfrac12\partial_{xx}(s^2m)\Bigr]dx = 0
  \qquad\text{for every }\varphi .
\]
\textbf{2d. Conclude.} The bracket must vanish everywhere:
\[
  \boxed{\;m_t = -\partial_x\bigl(b\,m\bigr) + \tfrac12\,\partial_{xx}\bigl(s^2m\bigr),
  \qquad m(0,x)=m_0(x).\;}
\]
The derivatives act on the products $bm$ and $s^2m$, not on $m$ alone. Both $b$ and $s^2$
depend on $x$ through $\pi^*$, which changes sharply near $\mb$, so this matters.

\textbf{Duality with the HJB.} For fixed $\pi$ the HJB uses
$\mathcal L V = bV_x + \tfrac12s^2V_{xx}$; the Fokker--Planck equation uses its
\emph{adjoint} $\mathcal L^*m = -\partial_x(bm)+\tfrac12\partial_{xx}(s^2m)$, defined by
$\int(\mathcal L\varphi)\,m = \int\varphi\,(\mathcal L^*m)$, which is exactly Step 2c.
$\mathcal L$ backward, $\mathcal L^*$ forward: this pairing is the structure of every
mean-field game.

\subsection*{Step 3: Conservation form}
\addcontentsline{toc}{subsection}{Step 3: Conservation form}
\[
  m_t + \partial_xJ = 0,\qquad J = b\,m - \tfrac12\,\partial_x\bigl(s^2m\bigr).
\]
$J(t,x)$ is the \textbf{flux}: investors per unit time crossing wealth level $x$ upward.
Density changes only through flow; nobody is created or destroyed.
\begin{itemize}
  \item \textbf{Drift part $bm$:} everyone is carried by expected growth.
  \item \textbf{Diffusion part $-\tfrac12\partial_x(s^2m)$:} randomness spreads investors from
        crowded to sparse regions, more strongly where $\pi^*$ is large.
\end{itemize}
\textbf{For $\lambda=0$:} $b\approx0.02x+0.09$ and $s^2\approx(1.5\times0.2)^2=0.09$ for
everyone, so the bell curve shifts right and widens evenly, staying Gaussian.

\textbf{With the penalty (a prediction to test):} investors just below $\mb$ hold little stock
(small $s^2$) and barely spread; those far behind hold more than Merton and spread more. Mass
tends to build up where the variance is small, so near $T$ the distribution should bunch up
just below the average and grow a fatter left tail.

\begin{caution}[The FP needs derivatives of the position]
$\partial_{xx}(\pi^{*2}\sigma^2m)$ involves $\partial_{xx}(\pi^{*2})$, and
$\pi^*\propto V_x/V_{xx}$. Expanded naively, the FP residual needs $V$ up to its
\emph{fourth} derivative, which is fragile for a PINN. The finite-volume flux form avoids
this (it only differences $s^2m$ at grid points). For a PINN, options are a separate network
for $\pi$, writing the FP loss in flux form, or computing the population side by Monte Carlo
and using the PINN only for $V$.
\end{caution}

\subsection*{Step 4: Initial and boundary conditions}
\addcontentsline{toc}{subsection}{Step 4: Initial and boundary conditions}
\textbf{Initial condition.} $m(0,x)=\mathcal N(1,0.2^2)$. The FP equation runs
\emph{forward} from it, opposite to the HJB.

\textbf{Boundaries: zero flux.} Use $J=0$ at $x=-1$ and $x=3$: nobody leaves the domain, and
$\frac{d}{dt}\int m\,dx = -[J]_{-1}^{3} = 0$ conserves mass exactly. (Setting $m=0$ instead
also works but leaks a little mass.) For $\lambda=0$ the true mass outside the domain at $T$
is tiny: about $3\times10^{-9}$ below $-1$ (5.8 standard deviations) and $10^{-7}$ above 3
(5.2 standard deviations). With the penalty the left tail is fatter, so monitor the mass in
the outermost cells; if it isn't negligible, widen the domain.

This choice is much easier than for the HJB, because nobody lives near the edges.

\subsection*{Step 5: Mean and variance, checked at \texorpdfstring{$\lambda=0$}{lambda = 0}}
\addcontentsline{toc}{subsection}{Step 5: Mean and variance, checked at lambda = 0}
Take $\varphi(x)=x$ in Step 2b (allowed because $m$ decays fast). Since $\varphi'=1$,
$\varphi''=0$:
\[
  \boxed{\;\frac{d\mb}{dt} = r\,\mb(t) + (\mu-r)\int\pi^*(t,x)\,m(t,x)\,dx .\;}
\]
This holds for \emph{any} $\lambda$: the average grows at the risk-free rate plus the risk
premium times the \emph{average position}. With the penalty, near-losers hold less and deep
losers hold more, so whether $\mb_T$ ends above or below 1.11 depends on the balance.

\textbf{At $\lambda=0$,} $\pi^*=1.5\,e^{-r(T-t)}$ for everyone, so
$\mb'=r\mb+0.09\,e^{-r(T-t)}$, giving
\[
  \mb(t) = e^{rt} + 0.09\,t\,e^{r(t-T)},\qquad \mb_T \approx 1.110 .
\]
Similarly $\varphi=x^2$ gives the variance $v(t)$: $v' = 2rv + \pi^{*2}\sigma^2
= 2rv + 0.09\,e^{-2r(T-t)}$, so
\[
  v(t) = 0.04\,e^{2rt} + 0.09\,t\,e^{2r(t-T)} .
\]
In fact, at $\lambda=0$ the exact solution is $m(t,\cdot)=\mathcal N(\mb(t),v(t))$, so an FP
solver can be checked pointwise, not just through moments. (\texttt{population\_mean} and
\texttt{population\_var} in \texttt{benchmarks.py}.)

\subsection*{Step 6: The coupled system and its fixed point}
\addcontentsline{toc}{subsection}{Step 6: The coupled system and its fixed point}
\label{sec:system}
\begin{keyidea}[The mean-field game system]
\vspace{-6pt}
\begin{align*}
  &\text{HJB (backward):} && V_t + rxV_x - \tfrac12\kappa^2\,V_x^2/V_{xx} = 0,
     && V(T,x) = U(x,\mb_T),\\
  &\text{position:}       && \pi^* = -(\mu-r)\,V_x/(\sigma^2V_{xx}), \\
  &\text{FP (forward):}   && m_t = -\partial_x\bigl([rx+\pi^*(\mu-r)]m\bigr)
                               + \tfrac12\partial_{xx}\bigl(\pi^{*2}\sigma^2m\bigr),
     && m(0,x) = m_0(x),\\
  &\text{consistency:}    && \mb_T = \textstyle\int x\,m(T,x)\,dx .
\end{align*}
\end{keyidea}

\textbf{The coupling is a single number.} Each investor's best response depends on the
population only through $\mb_T$. So define
\[
  \Phi(\mb) = \text{average terminal wealth when everyone best-responds to }\mb ,
\]
computed by: solve the HJB with $V(T,x)=U(x,\mb)$, read off $\pi^*$, solve the FP, take the
mean at $T$. An equilibrium is a solution of
\[
  \Phi(\mb) = \mb ,
\]
a \textbf{one-dimensional root-finding problem}. Options: damped iteration
$\mb_{k+1}=(1-\omega)\mb_k+\omega\,\Phi(\mb_k)$, or the secant method on $\Phi(\mb)-\mb$,
starting from the $\lambda=0$ value $1.110$. Plotting $\Phi(\mb)-\mb$ over a range of $\mb$
also shows whether the equilibrium is unique, which isn't guaranteed a priori.

By Step 5, $\Phi$ only needs the average position over time, so a Monte Carlo simulation of
investors under $\pi^*$ can stand in for the FP solver when computing $\Phi$.

\subsection*{Step 7: What it means for the solvers}
\addcontentsline{toc}{subsection}{Step 7: What it means for the solvers}
\begin{center}
\small
\begin{tabular}{@{}p{0.13\textwidth}p{0.47\textwidth}p{0.34\textwidth}@{}}
\toprule
Solver & Approach & Checks \\
\midrule
HJB, FD & implicit Euler backward in time; policy iteration (one tridiagonal solve per
iteration); boundaries of HJB Step 7 (details in Appendix~\ref{app:fd-hjb})
& matches \texttt{merton\_value} at $\lambda=0$; boundary-moving test \\
FP, FD & finite-volume (flux) form, zero flux at the edges; $\pi^*$ evaluated at cell
faces from the HJB solution
& mass conserved; Gaussian solution at $\lambda=0$ \\
Monte Carlo & simulate many investors under $\pi^*$ (interpolated)
& mean and variance at $\lambda=0$; matches FD density for $\lambda>0$ \\
Equilibrium & secant or damped iteration on $\Phi(\mb)=\mb$
& $\lambda=0$ gives $1.110$; $\Phi(\mb)-\mb$ plotted \\
HJB, PINN & network $V_\theta(t,x)$; loss = PDE residual + terminal + boundaries;
$\pi^*$ from automatic differentiation
& $\lambda=0$ closed form; FD for $\lambda>0$ \\
FP, PINN & network $m_\theta(t,x)$ on the FP residual
& Gaussian at $\lambda=0$; FD and Monte Carlo for $\lambda>0$ \\
\bottomrule
\end{tabular}
\end{center}

%======================================================================
\appendix
\section{The finite-difference HJB solver}
\label{app:fd-hjb}

This appendix turns the HJB of Section~2 into the scheme implemented in
\texttt{src/mfg/fd/hjb.py}. The terminal value $\mb$ is a fixed number throughout. Each step
says what the code does and why.

\subsection*{Step 1: The grid and the unknowns}
\addcontentsline{toc}{subsection}{Step 1: The grid and the unknowns}
Replace $(t,x)$ by grid points:
\[
  x_j = x_{\min} + j\,\Delta x,\quad j=0,\dots,J,\qquad
  t_0=0<t_1<\dots<t_N=T,\quad \Delta t_n = t_{n+1}-t_n ,
\]
\[
  V_j^n \approx V(t_n,x_j),\qquad \pi_j^n \approx \pi^*(t_n,x_j),
\]
stored as arrays of shape $(N+1,J+1)$. The time steps are uniform to begin with
($\Delta t_n=T/N$), but the scheme doesn't need that.

\textbf{Default resolution.} $J=800$, $N=400$: $\Delta x=0.005$, $\Delta t=0.0025$, about
$3\times10^5$ values, solved in seconds. With $x_{\min}=-1$ the kink $\mb=1.1$ is exactly the
grid point $j=420$. This isn't required, but it makes the discrete terminal condition
cleaner.

\textbf{Is it fine enough near the kink?} By HJB Step 9b the bump in $A_V$ has width about
$0.3\sqrt{T-t}$. One step before $T$ that is $0.3\sqrt{0.0025}=0.015$, about three cells, so
the last few steps are only coarsely resolved; everywhere else the resolution is
comfortable. The convergence study (Step 8) shows whether this matters.

\textbf{Known values on the edges of the grid.}
\begin{center}
\begin{tabular}{@{}lll@{}}
\toprule
Where & Value & Function \\
\midrule
top row, $t=T$ & $V_j^N = U(x_j,\mb)$, $\pi_j^N=\pi^*(T,x_j)$ & \texttt{utility}, \texttt{pi\_terminal} \\
left column, $x=-1$ & $V_0^n = V_{\text{fix}}(t_n,x_{\min})$ & \texttt{merton\_strategy\_value} \\
right column, $x=3$ & $V_J^n = V_{\text{CARA}}(t_n,x_{\max})$ & \texttt{merton\_value} \\
\bottomrule
\end{tabular}
\end{center}
Everything else is unknown.

\subsection*{Step 2: Time stepping}
\addcontentsline{toc}{subsection}{Step 2: Time stepping}
Write the HJB as $V_t + H(V)=0$ with
$H(V)=\sup_\pi\{[rx+\pi(\mu-r)]V_x+\tfrac12\pi^2\sigma^2V_{xx}\}$. We know row $n+1$ and want
row $n$ (backward in time). Replacing $V_t$ by $(V^{n+1}-V^n)/\Delta t_n$, the choice is
where to evaluate $H$:
\begin{center}
\small
\begin{tabular}{@{}lp{0.34\textwidth}p{0.36\textwidth}@{}}
\toprule
Scheme & Update & Here \\
\midrule
Explicit & $V^n = V^{n+1} + \Delta t\,H(V^{n+1})$
  & stable only if $\Delta t\le\Delta x^2/s^2_{\max}\approx1.3\times10^{-4}$: about 7700
    steps, and fragile \\
\textbf{Implicit Euler} & $V^n - \Delta t\,H(V^n) = V^{n+1}$
  & stable for any $\Delta t$; \textbf{used} \\
Crank--Nicolson & average of the two
  & second order, but oscillates near the kink \\
\bottomrule
\end{tabular}
\end{center}
Reasons for implicit Euler:
\begin{itemize}
  \item \textbf{Stable for any step,} so $\Delta t$ is chosen for accuracy, not to avoid
        blow-up.
  \item \textbf{Monotone} (with the space derivatives of Step 4): it creates no wiggles.
        $\pi^*$ is read off $V_{xx}$, so wiggles in $V$ would become wild positions.
  \item \textbf{Safe with the kink.} Crank--Nicolson is known to oscillate when the terminal
        data has a kink, and the usual cure is a few implicit Euler steps first anyway.
  \item The price is first-order accuracy in time. If the convergence study shows the time
        error dominates, use smaller steps near $T$, where $V$ changes fastest.
\end{itemize}
Each time step therefore solves, at every interior node $j$,
\[
  \boxed{\;V_j^n - \Delta t_n\sup_{\pi}\Bigl\{b_j(\pi)\,(D_xV^n)_j
      + \tfrac12 s_j^2(\pi)\,(D_{xx}V^n)_j\Bigr\} = V_j^{n+1},\;}
\]
with $b_j(\pi)=rx_j+\pi(\mu-r)$, $s_j^2(\pi)=\pi^2\sigma^2$, and $D_x$, $D_{xx}$ the difference
operators of Step 4 applied to the \emph{unknown} row $V^n$. Because of the sup, this is
nonlinear in $V^n$.

\subsection*{Step 3: Policy iteration inside each time step}
\addcontentsline{toc}{subsection}{Step 3: Policy iteration inside each time step}
For a \emph{fixed} policy the problem is linear, so alternate two easy steps.

\textbf{(a) Fix $\pi$, solve for $V$.} With $\pi^{(k)}$ given, $b_j$ and $s_j^2$ are numbers,
and
\[
  (L^{(k)}V)_j = b_j\,(D_xV)_j + \tfrac12 s_j^2\,(D_{xx}V)_j
\]
is a linear operator. Then $V^{(k)}$ solves $(I-\Delta t_n L^{(k)})\,V^{(k)} = V^{n+1}$: one
tridiagonal solve (Step 5).

\textbf{(b) Fix $V$, update $\pi$.} The sup is a separate parabola in $\pi$ at each node, so
\[
  \pi_j^{(k+1)} = -\frac{(\mu-r)\,(D_xV^{(k)})_j}{\sigma^2\,(D_{xx}V^{(k)})_j},
\]
the discrete version of HJB Step 5 (with the guard of Step 6). It is one array expression.

\begin{keyidea}[One time step]
\begin{enumerate}[label=\arabic*.,leftmargin=*]
  \item Warm start: $\pi \leftarrow \pi^{n+1}$ (at the first step, $\pi^N$ from
        \texttt{pi\_terminal}, which already jumps from 1.5 to 0.214 at $\mb$).
  \item Solve $(I-\Delta t_n L^{\pi})\,V = V^{n+1}$ with the boundary values at $t_n$.
  \item Update $\pi$ from $D_xV$ and $D_{xx}V$.
  \item Stop if $\max_j|V_j-V_j^{\text{old}}|<\text{tol}$ or after \texttt{max\_iter}
        iterations; otherwise go to 2.
  \item Store $V^n\leftarrow V$, $\pi^n\leftarrow\pi$ and the iteration count.
\end{enumerate}
\end{keyidea}

\textbf{Why it converges.}
\begin{enumerate}[label=(\roman*)]
  \item \textbf{$V$ only goes up.} The new policy maximises the bracket at every node, so
        $L^{(k+1)}V^{(k)}\ge L^{(k)}V^{(k)}$, hence
        $(I-\Delta tL^{(k+1)})V^{(k)} \le V^{n+1} = (I-\Delta tL^{(k+1)})V^{(k+1)}$. If
        $I-\Delta tL$ is an \textbf{M-matrix} (nonnegative inverse; Step 4), the inverse
        keeps the inequality: $V^{(k+1)}\ge V^{(k)}$. A better strategy never lowers the value.
        $V$ is also bounded above (by the CARA value), so the iteration can't run away.
  \item \textbf{It is Newton's method in disguise.} With
        $F(V)=V-\Delta t\sup_\pi L^\pi V - V^{n+1}$, the envelope theorem says the derivative of
        the sup is the operator at the maximising $\pi$, so step (a) is a Newton step on
        $F=0$. Convergence is fast near the solution, and the warm start puts us there; the
        code records the iteration counts so this can be checked.
\end{enumerate}

\textbf{Lagged policy.} Stopping after one solve (\texttt{max\_iter=1}) uses $\pi^{n+1}$ for
step $n$. It is still consistent and first order in time, but lags where $\pi$ moves fast
(near the kink, close to $T$). The test of Step 8 checks that both versions agree as
$\Delta t\to0$.

\subsection*{Step 4: Space derivatives and the M-matrix check}
\addcontentsline{toc}{subsection}{Step 4: Space derivatives and the M-matrix check}
At an interior node, the second derivative is always central:
\[
  (D_{xx}V)_j = \frac{V_{j+1}-2V_j+V_{j-1}}{\Delta x^2}.
\]
For the first derivative there are two choices:
\[
  \text{central: } \frac{V_{j+1}-V_{j-1}}{2\Delta x}\ \ (\text{error }O(\Delta x^2)),
  \qquad
  \text{upwind: } \begin{cases}\dfrac{V_{j+1}-V_j}{\Delta x}, & b_j>0,\\[8pt]
                                \dfrac{V_j-V_{j-1}}{\Delta x}, & b_j<0,\end{cases}
  \ \ (\text{error }O(\Delta x)).
\]
(Upwind looks toward larger $x$ when $b>0$ because the HJB runs backward: the value at $x$
depends on where the investor drifts \emph{to}.)

\textbf{Coefficients.} Either way, $L$ couples each node to its two neighbours:
\[
  (LV)_j = \alpha_j\,V_{j-1} - (\alpha_j+\beta_j)\,V_j + \beta_j\,V_{j+1},
\]
\begin{center}
\begin{tabular}{@{}lll@{}}
\toprule
 & $\alpha_j$ (to $j-1$) & $\beta_j$ (to $j+1$) \\
\midrule
central & $\dfrac{s_j^2}{2\Delta x^2}-\dfrac{b_j}{2\Delta x}$
        & $\dfrac{s_j^2}{2\Delta x^2}+\dfrac{b_j}{2\Delta x}$ \\[10pt]
upwind  & $\dfrac{s_j^2}{2\Delta x^2}+\dfrac{\max(-b_j,0)}{\Delta x}$
        & $\dfrac{s_j^2}{2\Delta x^2}+\dfrac{\max(b_j,0)}{\Delta x}$ \\
\bottomrule
\end{tabular}
\end{center}
The rows of $L$ sum to zero (a constant has zero derivatives).

\textbf{The M-matrix condition.} The matrix $A=I-\Delta tL$ has diagonal
$1+\Delta t(\alpha_j+\beta_j)$ and off-diagonals $-\Delta t\alpha_j$, $-\Delta t\beta_j$.
If $\alpha_j,\beta_j\ge0$, the off-diagonals are $\le0$ and every row sums to $1>0$, which
makes $A$ an M-matrix: $A^{-1}\ge0$ entrywise. Two consequences:
\begin{itemize}
  \item the policy-iteration argument of Step 3(i) holds;
  \item the scheme is \textbf{monotone}: a larger $V^{n+1}$ gives a larger $V^n$, so there
        are no new wiggles and no new maxima.
\end{itemize}
Upwind always has $\alpha_j,\beta_j\ge0$. Central needs the \textbf{P\'eclet condition}
\[
  s_j^2 \ge |b_j|\,\Delta x \quad\Longleftrightarrow\quad \text{diffusion beats drift on the
  scale of one cell.}
\]
\textbf{Rule used:} central where the P\'eclet condition holds (second order), upwind where
it fails (always monotone). With $\Delta x=0.005$, central works wherever $\pi\gtrsim0.06$
near $\mb$ (at worst $\pi\gtrsim0.09$, at $x=3$). The smallest position we expect is 0.214,
so central should be used everywhere and upwind is only a safety net; the code counts the
upwind nodes, and the tests expect zero.

\subsection*{Step 5: The tridiagonal system}
\addcontentsline{toc}{subsection}{Step 5: The tridiagonal system}
Each solve in Step 3(a) is $AV=\text{rhs}$ over all $J+1$ nodes. Writing
$a_j=\Delta t_n\alpha_j$ and $c_j=\Delta t_n\beta_j$:
\[
\begin{pmatrix}
1 & 0 \\
-a_1 & 1+a_1+c_1 & -c_1 \\
 & \ddots & \ddots & \ddots \\
 & & -a_{J-1} & 1+a_{J-1}+c_{J-1} & -c_{J-1} \\
 & & & 0 & 1
\end{pmatrix}
\begin{pmatrix} V_0 \\ V_1 \\ \vdots \\ V_{J-1} \\ V_J \end{pmatrix}
=
\begin{pmatrix} V_{\text{fix}}(t_n,x_{\min}) \\ V_1^{n+1} \\ \vdots \\ V_{J-1}^{n+1} \\
                V_{\text{CARA}}(t_n,x_{\max}) \end{pmatrix}.
\]
\begin{itemize}
  \item \textbf{Dirichlet rows.} The first and last rows just say ``$V$ equals the boundary
        value at $t_n$''. They keep the M-matrix property.
  \item \textbf{Tridiagonal.} Each node talks only to its neighbours, so the solve costs
        $O(J)$ (Thomas algorithm), via \texttt{scipy.linalg.solve\_banded} with the three
        diagonals stacked in an array of shape $(3,J+1)$.
  \item \textbf{Cost.} $N$ steps $\times$ (a few iterations) $\times$ $O(J)$: well under a
        second per iteration sweep at the default grid.
\end{itemize}

\subsection*{Step 6: Reading off the position, with a guard}
\addcontentsline{toc}{subsection}{Step 6: Reading off the position, with a guard}
At the interior nodes, using the \emph{central} $D_x$:
\[
  \pi_j = \operatorname{clip}\Bigl(-\frac{(\mu-r)(D_xV)_j}{\sigma^2(D_{xx}V)_j},\ 0,\ \pi_{\max}\Bigr)
  \quad\text{if }(D_{xx}V)_j<0,\qquad \pi_j=\pi_{\max}\ \text{otherwise.}
\]
\begin{itemize}
  \item \textbf{Why a box $[0,\pi_{\max}]$.} If the discrete $V$ were locally convex, the
        sup would be $+\infty$. In theory it can't be: $U$ is strictly concave (only $U''$
        jumps at $\mb$, by $-\lambda$) and the HJB preserves concavity (Step 9d). On a
        bounded box the maximum always exists: the clipped formula when $D_{xx}V<0$, the
        upper edge otherwise (where $V_x>0$). Default $\pi_{\max}=10$, far above the
        $\approx2.2$ expected at $x=-1$.
  \item \textbf{Clipping is a symptom.} The code counts clipped nodes; the tests expect zero.
  \item \textbf{Edge nodes.} $\pi$ isn't needed at $j=0,J$ (Dirichlet rows), but the
        Fokker--Planck solver wants it everywhere, so it is copied from the neighbouring node.
        Positions at the edges are unreliable anyway (see the boundary-moving test).
  \item \textbf{At $t=T$,} $\pi^N$ comes from \texttt{pi\_terminal} exactly, not from
        differencing the kinked $U$.
\end{itemize}

\subsection*{Step 7: The solution object}
\addcontentsline{toc}{subsection}{Step 7: The solution object}
Following \texttt{CONTRIBUTING.md}, the solver returns an object the PINN side can be
compared with:
\begin{center}
\small
\begin{tabular}{@{}lp{0.62\textwidth}@{}}
\toprule
Item & Meaning \\
\midrule
\texttt{solve\_hjb(p, mbar, nx, nt, ...)} & runs Steps 1--6; returns the object below \\
\texttt{t}, \texttt{x} & the grids, shapes $(N+1)$ and $(J+1)$ \\
\texttt{V}, \texttt{pi} & the arrays $V_j^n$, $\pi_j^n$, shape $(N+1,J+1)$ \\
\texttt{value(t, x)}, \texttt{position(t, x)} & bilinear interpolation of \texttt{V}, \texttt{pi};
  NumPy arrays broadcasting like \texttt{t + x} \\
\texttt{iterations} & policy iterations per time step \\
\texttt{n\_clipped}, \texttt{n\_upwind} & guard counters of Steps 4 and 6 (expected 0) \\
\bottomrule
\end{tabular}
\end{center}
Linear interpolation in $x$ has error $O(\Delta x^2)$, the same order as the scheme.

\subsection*{Step 8: Tests and convergence study}
\addcontentsline{toc}{subsection}{Step 8: Tests and convergence study}
\begin{center}
\small
\begin{tabular}{@{}p{0.27\textwidth}p{0.66\textwidth}@{}}
\toprule
Test & What it checks \\
\midrule
$\lambda=0$: $V$ and $\pi$ & against \texttt{merton\_value} and \texttt{merton\_pi}: the whole
  scheme on a problem with a known answer \\
Convergence at $\lambda=0$ & refine $(J,N)=(200,100),(400,200),(800,400)$: the error should
  roughly halve each time (implicit Euler is first order in time, which dominates) \\
Sandwich at $\lambda=2$ & $V_{\text{fix}}\le V\le V_{\text{CARA}}$ everywhere: no strategy
  beats optimal, and the penalty only lowers utility \\
Shape at $\lambda=2$ & $D_xV>0$ and $D_{xx}V<0$ at every node; no clipping, no upwinding \\
Near $T$ at $\lambda=2$ & one step before $T$, $\pi$ matches \texttt{pi\_terminal} away from
  the kink and the edge node \\
Policy vs lagged & full policy iteration and \texttt{max\_iter=1}: $V$ agrees; $\pi$ agrees for
  $t\le0.9$ \\
Boundary-moving & $x_{\min}=-1$ vs $-2$: $V$ and $\pi$ agree on $[0,2.2]$, where the
  population lives \\
\bottomrule
\end{tabular}
\end{center}
At $\lambda=2$ there is no closed form, so a refinement study (same three grids, differences
between successive solutions) gives the accuracy of the reference solution the PINN will be
compared against.

\textbf{What the first runs showed} (default grid $J=800$, $N=400$; about 0.2\,s):
\begin{itemize}
  \item \textbf{$\lambda=0$:} largest errors $2.4\times10^{-6}$ in $V$ and $3.3\times10^{-4}$ in
        $\pi$. Refining $(400,200)\to(800,400)\to(1600,800)$ divides them by 1.7--1.9: first
        order, as expected.
  \item \textbf{$\lambda=2$:} three to four policy iterations per step (up to six at the first
        step below $T$), no clipping, no upwinding, and $V$ concave and increasing at every
        node. $V(0,1.1)$ changes by less than $10^{-6}$ between the two finest grids.
  \item \textbf{Positions at $t=0$:} $\pi=1.86$ at $x=-0.5$, $1.00$ at $x=0.5$, $0.65$ at the
        average $x=1.1$ and $1.47$ (Merton) at $x=2$. The fear zone is wide at $t=0$: an
        investor exactly at the average holds less than half the Merton position.
  \item \textbf{The edge node.} Near $x=-1$ the solver's $\pi$ is pulled towards the boundary
        strategy (Merton, 1.5) instead of the optimal 2.2. On $[0,2.2]$ moving the boundary to
        $-2$ changes $V$ by $2.5\times10^{-5}$ and $\pi$ by $1.5\times10^{-3}$, but at
        $x=-0.5$ it changes $\pi$ by 0.04: at $t=0$ the boundary's influence reaches about half
        a unit into the domain.
  \item \textbf{Lagged policy.} $V$ agrees with full policy iteration to $10^{-5}$, and $\pi$
        to $10^{-3}$ for $t\le0.9$. At the kink in the first steps below $T$ the two differ by
        $O(1)$ in $\pi$: the lagged step uses the terminal policy, which jumps there. Full
        policy iteration is the default.
\end{itemize}

\end{document}
